\chapter{字符串与格式化}
\label{ch:12}

字符串既是文本，也是一段需要管理长度和容量的数据。构造页面、解析文件和显示错误之前，需要先区分字符语义、字节边界与格式化协议。

\section{文本与边界}

String 使用 UTF-8。处理切片或外部字节长度时注意多字节编码：任意字节偏移未必是字符边界。输入长度限制也要明确采用字节还是用户可见字符；不要仅凭 ASCII 例程推断中文和表情的行为。

\section{插值与转义}

格式串的占位符由参数填入，内嵌表达式按语言语法求值。需要输出花括号本身时使用对应转义。格式化便于展示，但不是 HTML 转义、SQL 引号处理或 JSON 序列化的替代品。

\par
\medskip
\begin{lstlisting}[style=chali]
void main()
{
    int64 count = 3;
    String text = f"count={count + 1}, {f"inner={count}"}, \{literal\}";
    println("{}", text);
}
\end{lstlisting}
\par
\medskip

\par
\medskip
\begin{lstlisting}[style=out]
count=4, inner=3, {literal}
\end{lstlisting}
\par
\medskip

\section{结构化显示与自定义显示}

默认结构化显示适合诊断，文本形式不应当作稳定的文件格式。需要业务显示时提供 toString；需要持久化或跨进程传输时使用相应序列化协议并处理失败。

\par
\medskip
\begin{lstlisting}[style=chali]
struct Pair
{
    String left;
    String right;
};

void main()
{
    Pair pair;
    pair.left = "one";
    pair.right = "two";
    List<Any> list = [1, "two", true];
    Set<int64> set = {3, 1, 2};
    Map<String, Any> map = {"name": "ChaLi", "count": 3};
    let tuple = (1, "two", list);

    println("{}", pair.toString());
    println("{}", list.toString());
    println("{}", set.toString());
    println("{}", map.toString());
    println("{}", tuple.toString());
}
\end{lstlisting}
\par
\medskip

\par
\medskip
\begin{lstlisting}[style=out]
{"one", "two"}
[1, "two", true]
{3, 1, 2}
{"name": "ChaLi", "count": 3}
(1, "two", [1, "two", true])
\end{lstlisting}
\par
\medskip

\par
\medskip
\begin{lstlisting}[style=chali]
struct Named
{
    String name;
    String toString() { return "name=" + self.name; }
};

void main()
{
    Named value;
    value.name = "ChaLi";
    String text = value;
    println("{}", text);
    println("{}", cast<String>(value));
}
\end{lstlisting}
\par
\medskip

\par
\medskip
\begin{lstlisting}[style=out]
name=ChaLi
name=ChaLi
\end{lstlisting}
\par
\medskip

\section{拼接、缓冲区与安全输出}

重复拼接的成本取决于复制、临时值和容量增长。String 可以复用已分配的缓冲区，缩短内容不等于释放容量；优化应保持复制和借用的可观察行为。第 24 章区分编译选项，第 25 章网页示例通过 textContent 显示用户数据，避免将姓名直接解释成 HTML。

\section{习题}

\begin{enumerate}

\item 输出一个包含中文的错误消息。

\item 格式化是否能安全拼 SQL？

\item 何时使用 toString？

\item 缩短字符串后需要立即释放缓冲区吗？

\end{enumerate}

\subsection*{参考解答与验收要点}

\begin{enumerate}

\item 检查终端编码和实际输出，不按字节数猜字符数。

\item 不能；文本值必须按数据库接口转义或参数绑定。

\item 用于显示；稳定存储使用明确的编码格式。

\item 不需要；保留容量可供后续增长，drop 清理最终缓冲区。

\end{enumerate}

下一章：\hyperref[ch:13]{第 13 章}。

\section{番外：施氏数“字符”}

借赵元任《施氏食狮史》的同音构思另写，不复述原文。

施氏试字符串，示十式，使实习生识之。生视“施施施”，说是三；施氏试 \texttt{len}，示九。生失色，疑机器失识。施氏说：“是三个汉字，实是九字节。”

生始释，遂试切。使首字节至第三字节，各自成串。屏上似石似齿，实皆失形。生说：“既许切，何事失真？”施氏说：“许切是字节接口的事，切得像不像字，是编码的事。”

生再试十个字母，十个字节；再试十个汉字，长度又变。于是施氏施一纸示例：字节非码点，码点非字形。生读三次，仍将“字符数”填进“存储大小”。施氏只好把表格改成三个字段，免得同音字一起进生产。

此后实习生每逢产品说“就限制十个字，十分简单”，便失神片刻，先问：“十字节、十码点，还是十个看起来像一个的东西？”

产品说他话太多。施氏说：“是。谁先把单位写清楚，谁就少开十次会。”

\subsection*{技术对照}

String 的长度、下标和切片按当前 API 规定的单位使用。本章基础 String 操作按字节；处理 UTF-8 时不能任意截断多字节编码，也不能把代码点直接等同于视觉字形。
